\lesson{6}{Oct 5 2021 Tue (05:18:32)}{Narrative Writing}{Unit 1}

\subsubsection*{Elements of Narrative Writing}

Every good narrative uses these six traits to tell a story:

\begin{enumerate}
    \item \textbf{Ideas}: Key elements of the story and the details that support them
    \item \textbf{Organization}: Format, structure, and timeline of the story
    \item \textbf{Voice}: The way words and phrases are used to tell the story
    \item \textbf{Word Choice}: Descriptive words and phrases make the writing interesting
    \item \textbf{Sentence Fluency}: The way sentences are written and flow together to tell the story
    \item \textbf{Conventions}: The use of proper spelling, punctuation, capitalization, and grammar
\end{enumerate}

\subsubsection*{Ideas}

Ideas are the building blocks of a narrative, and they need to be developed throughout the story by supporting details.

\subsubsection*{Organization}

Organization takes all the details and puts them together so they make sense. Organization lets a reader know what's happening and when it
happens. Traditionally, stories have a clear beginning, middle, and end.

\subsubsection*{Voice and Word Choice}

Depending on where you are going, a writer has to choose the right words, tone, and mood to develop a voice.

\subsubsection*{Sentence Fluency and Conventions}

The word fluency sounds a lot like the word fluid for a good reason—sentence fluency helps a writer's ideas flow from beginning to end.
A good flow creates a rhythm that keeps a story interesting.

Proper use of conventions—spelling, grammar, capitalization, and punctuation—is an important part of creating sentence fluency.
Without proper conventions, a story is hard to read and difficult to understand.


\subsubsection*{The Writing Process}

\paragraph{Pre-write}

During the pre-writing process, you generate ideas, determine the voice needed to give life to your ideas, and organize the ideas to put them together to make a complete narrative.

\paragraph{Draft}

When you sit down to write the first draft of your story, you use all the ideas you generated during pre-writing. Sentence fluency is important in this step because you are figuring out how to weave your ideas together into a complete story.

\paragraph{Revise/Edit}

After you have finished your first draft, it is time to revise and edit—these are not the same thing. When you revise, you go back and look over what you've written to see how well you've done building your idea into a complete story. Does the voice you created match the story you are telling? Have you organized your writing so it makes sense? This is also a great time to share your story with someone else to get feedback about your writing.

When you edit, you check your writing for proper use of conventions—spelling, capitalization, grammar, and punctuation.

\paragraph{Rewrite}

In this stage of the writing process, you make any improvements needed based on your revision and editing. Then you check one last time to make sure your paper is free from conventions errors—and make corrections as needed.

\paragraph{Publish}

Your hard work and attention to detail have paid off. Now is the time to send your story out into the world for others to read.

\newpage
